\titledquestion{Stacks and Queues Using Linked Lists}

Use acyclic singly linked lists. For each list, retain only its head pointer; do not keep a tail pointer. You may use temporary node pointers. Explain your algorithms in pseudocode or clear English. Handle an empty structure explicitly.

\begin{parts}
\part[4] Implement a stack using one such list. Describe \ttt{push(x)} and \ttt{pop()}, where \ttt{pop()} removes and returns the top value. Each operation must take constant time.
\begin{solution}

\begin{lstlisting}[language=C++]
template <typename T> class Stack {
private:
struct Node {
    T data;
    Node *nxt;
    Node(const T &_data, Node *_nxt = nullptr) : data(_data), nxt(_nxt) {}
};

// Keep top of the stack
Node *head;

public:
Stack() : head(nullptr) {}
~Stack() {
    while (!isEmpty()) {
    pop();
    }
}

bool isEmpty() const { return head == nullptr; }

void push(const T &value) { head = new Node(value, head); }

T pop() {
    if (isEmpty())
    throw std::runtime_error("Stack is empty");

    Node *temp = head;
    T pop_data = head->data;
    head = head->nxt;
    delete temp;

    return pop_data;
}

T top() const {
    if (isEmpty())
    throw std::runtime_error("Stack is empty");

    return head->data;
}
};
\end{lstlisting}
\end{solution}

\part[4] Implement a queue using two stacks from part (a), called \ttt{in} and \ttt{out}. Describe \ttt{enqueue(x)} and \ttt{dequeue()}, and explain why values leave in FIFO order. Move values between the stacks only when necessary.
\begin{solution}

\begin{lstlisting}[language=C++]
template <typename T> class Queue {
private:
    Stack<T> in, out;

public:
    bool isEmpty() const { return in.isEmpty() && out.isEmpty(); }

    void enqueue(const T &value) { in.push(value); }

    T dequeue() {
        if (out.isEmpty()) {
        while (!in.isEmpty()) {
            out.push(in.pop());
        }
        }
        if (out.isEmpty())
        throw std::runtime_error("Queue is empty");
        return out.pop();
    }
};
\end{lstlisting}

enqueue(x) pushes x onto the \ttt{in} stack.

dequeue() pops from the \ttt{out} stack. If \ttt{out} is empty, it transfers all elements from \ttt{in} to \ttt{out}, reversing their order. This ensures FIFO behavior, as the first element enqueued will be at the top of \ttt{out} after the transfer.

Stack is a LIFO structure, so we inplement the LIFO behavior twice to reverse the order, thus achieving FIFO structure.
\end{solution}

\part[2] Starting with an empty queue, perform
\[
\mathtt{enqueue(1),\ enqueue(2),\ dequeue(),\ enqueue(3),\ dequeue().}
\]
Give the two returned values and the final contents of each stack, from bottom to top.
\begin{solution}
\begin{enumerate}
  \item[(1)] enqueue(1): in = [1], out = []
  \item[(2)] enqueue(2): in = [1, 2], out = []
  \item[(3)] dequeue(): out is empty, transfer from in to out:
  \begin{enumerate}
    \item[] pop 2 from in, push to out: in = [1], out = [2]
    \item[] pop 1 from in, push to out: in = [], out = [2, 1]
    \item[] pop from out: return 1, now in = [], out = [2]
  \end{enumerate}
  \item[(4)] enqueue(3): in = [3], out = [2]
  \item[(5)] dequeue(): pop from out: return 2, now in = [3], out = []

\end{enumerate}
\end{solution}
\end{parts}
